\section{结构基础与多晶型}

\subsection{低温相与高温相}

BaZn$_2$Si$_2$O$_7$ 在约 280--300~$^\circ$C 发生结构转变；低温相的 ZnO$_4$ 链具有较强的方向性，高温相则通过四面体转动和链间角度调整释放热应变 \cite{lin1999phase,thieme2015ba1}。HT\nobreakdash-XRD 与热膨胀测量表明，两种结构的热膨胀张量差异远大于常规烧结误差，说明相变本身是相控制的首要变量 \cite{kerstan2012thermal,thieme2016very}。

公开结构数据将低温 BaZn$_2$Si$_2$O$_7$ 归入正交晶系 Cmc2$_1$ 空间群，其中 Ba$^{2+}$ 为五配位，Zn$^{2+}$ 和 Si$^{4+}$ 均为四面体配位；ZnO$_4$ 与 SiO$_4$ 通过角共享形成三维网络 \cite{available2020materials}。这些配位数据适合作为结构类比的第一层筛选依据，但数据库条目并非高温原位结构证据，不能替代原始衍射精修。表~\ref{tab:structure}\CJKecglue{}汇总了这些相的结构信息与可用于相控制的结论。

\begin{table}[bp]
\centering
\caption{BaZn$_2$Si$_2$O$_7$ 及其结构相关化合物的结构与热行为。}
\label{tab:structure}
\begin{tabular}{P{\dimexpr 0.2102\textwidth-2\tabcolsep\relax}P{\dimexpr 0.2985\textwidth-2\tabcolsep\relax}P{\dimexpr 0.1978\textwidth-2\tabcolsep\relax}P{\dimexpr 0.2935\textwidth-2\tabcolsep\relax}}
\toprule
    \thead{相或组成} & \thead{结构/配位信息} & \thead{主要表征} & \thead{可用于相控制的结论}\\
\midrule
    BaZn$_2$Si$_2$O$_7$ & 正交；Ba 五配位；ZnO$_4$\slash SiO$_4$ 四面体 & HT\nobreakdash-XRD、热膨胀测量 & 存在低温--高温转变 \cite{lin1999phase,available2020materials}\\
    \addlinespace[3pt]
    Ba$_{1-x}$Sr$_x$Zn$_2$Si$_2$O$_7$ & Sr 取代 Ba；高温结构可稳定至室温 & 单晶 XRD、HT\nobreakdash-XRD & $x\gtrsim0.1$ 时转变温度显著降低 \cite{thieme2015ba1,thieme2016very}\\
    \addlinespace[3pt]
    BaZn$_{2-x}$Mg$_x$Si$_2$O$_7$ & Mg 取代 Zn；结构随 $x$ 变化 & HT\nobreakdash-XRD、热膨胀测量 & 取代改变相变温度和热膨胀各向异性 \cite{kerstan2012thermal}\\
\bottomrule
\end{tabular}
\end{table}

\subsection{结构判据}

用于比较 Ba$_5$Y$_{12}$Zn[O(SiO$_4$)]$_8$ 的判据可分为四层：第一，硅酸根是否保持孤立四面体；第二，Zn 是否保持四配位并形成可追踪的角共享链；第三，Ba\slash Y 多面体是否具有相近的连接维度；第四，空间群和晶胞参数是否相近。前两层来自 BaZn$_2$Si$_2$O$_7$ 的直接结构证据 \cite{lin1999phase,available2020materials}，后两层需要 Ba$_5$Y$_{12}$Zn[O(SiO$_4$)]$_8$ 的单晶或高分辨粉末衍射数据，目前尚未获得。
